\begin{abstract}
Music understanding systems often combine text and audio, and audio is often represented as
a graph: short segments become nodes, with edges linking segments close in time or similar
in content. When a graph's structure is a deterministic function of its own node features,
the data-processing inequality guarantees it cannot add Shannon information beyond those
features -- an analytic fact, not something an experiment can test. What experiments can
test is the practical question: whether such a graph still helps a real, finite model, e.g.\
as a computational shortcut or inductive bias. We test this directly across four tasks
sharing the same node features and graph construction rule: text-only tag prediction,
graph-based genre classification, graph-plus-text fusion, and audio-text contrastive
retrieval. Tasks 2 and 4 each compare their graph-based model against a width/depth-matched,
lower-parameter edge-free alternative on held-out test data; Task 3's own comparison is
validation-only, so its evidence comes from two test-set diagnostics instead:
branch-zeroing inside the trained fusion model, and a cross-tabulation of
independently-trained single-branch checkpoints. No task shows a graph-based model with a
\emph{learned} aggregation step clearly ahead of its edge-free counterpart -- though a
non-learned exact 2-hop aggregation in genre classification outperforms both, using a
different training recipe than what it is compared against, so this suggests
under-exploited topology, not a clean ablation (reported, not hidden). In fusion, zeroing
the graph branch changes only 1.7\% of individual predictions and costs 0.009 macro-F1 --
small, not zero. A complementarity check first appears to show the graph-only model
correcting 48.5\% of the text-only model's errors, but splitting by label polarity shows it
predicts a tag present at all for only 5 of 50 tags, and fewer than 1\% of its apparent
corrections are genuine detections -- not real complementary signal as originally framed; a
follow-up check in a second implementation did not confirm this reading, discussed rather
than resolved. In retrieval, graph-based and no-graph encoders show no reliable difference
across three seeds, on every metric. We also checked two alternative explanations rather
than assume them away: caption text leaking target labels directly (it does,
substantially), and duplicate audio between train and test in one dataset (real, confirmed
for a small number of tracks, left open as an unresolved caveat for the smallest
comparison). We do not conclude graphs are useless for music -- only that this specific
construction does not reliably help beyond its own node features, for these tasks, on
these datasets. We report this as the paper's central finding, not a failed architecture
search.
\end{abstract}
